% Generated draft: requires mathematical and source review.
\documentclass[11pt]{article}
\usepackage[margin=25mm]{geometry}
\usepackage{fontspec}
\setmainfont{Latin Modern Roman}
\usepackage{amsmath,amssymb,mathtools,unicode-math}
\setmathfont{Latin Modern Math}
\usepackage{enumitem,needspace,xurl,hyperref}
\setlength{\parindent}{0pt}
\setlength{\parskip}{5pt}
\setlength{\emergencystretch}{3em}
\newcommand{\N}{\mathbb{N}}
\newcommand{\Z}{\mathbb{Z}}
\newcommand{\Q}{\mathbb{Q}}
\newcommand{\R}{\mathbb{R}}
\newcommand{\C}{\mathbb{C}}
\newcommand{\F}{\mathbb{F}}
\newcommand{\A}{\mathbb{A}}
\newcommand{\PP}{\mathbb{P}}
\newcommand{\E}{\mathbb{E}}
\newcommand{\eps}{\varepsilon}
\newcommand{\dd}{\,\mathrm{d}}
\begin{document}
\section*{UnsolvedMath ID: 30006033}
\textbf{Source problem code:} catalogue entry
\subsection{Definitions and mathematical statement}
The supplied record names two higher-categorical structures: fair \(n\)-categories and weakly globular \(n\)-fold categories. It does not supply their definitions, the admissible range of the dimension parameter \(n\), the comparison morphisms, or the intended notion of equivalence. Accordingly, let \(\operatorname{Fair}_n\) and \(\operatorname{WGlob}_n\) be schematic names for the respective theories; this notation does not specify categories or homotopy theories.

The proposed formulation of Simpson's weak-units conjecture asks whether, for the dimensions intended by the source, there is an equivalence schematically expressed as
\[
\operatorname{Fair}_n \simeq \operatorname{WGlob}_n.
\]
Here \(\simeq\) records the source's unspecified comparison: it must not be read as asserting an ordinary equivalence of categories, a model-categorical equivalence, or any other particular comparison without further source evidence. The supplied passage presents this as one possible formulation of Simpson's conjecture. A precise standalone conjecture requires definitions of both structures, the range and quantification of \(n\), and the meaning of equivalence.
\subsection{Short English statement}
The question asks whether fair \(n\)-categories and weakly globular \(n\)-fold categories give equivalent higher-categorical theories, as a proposed formulation of Simpson's weak-units conjecture. The supplied record leaves the definitions, dimension range, and meaning of equivalence unresolved.
\subsection{Sources}
\begin{itemize}
\item Local UnsolvedMath registry: supplied source record, permanent ID \(30006033\), title ``Simpson's Weak-Units Conjecture for Higher Categories,'' exact source problem code \texttt{catalogue entry}. The registry labels the problem open and its statement \texttt{corrected\_verified}; these are supplied curation labels, not independent mathematical verification.
\item Supplied report citation: \emph{Homotopical Algebra and Higher Structures (2025)}, Oberwolfach Reports, DOI 10.4171/owr/2024/39, pp.\ 2251--2308. \url{https://doi.org/10.4171/owr/2024/39}. The literature-source list labels this same reference ``Oberwolfach Report 39 (2024).'' The supplied passage describes the equivalence as one way to formulate Simpson's conjecture; no individual proposer of this extracted question is identified.
\item Supplied literature reference: Paoli, \emph{Weakly globular N-fold categories}. \url{https://arxiv.org/abs/1609.04072}.
\item Supplied literature assessment dated exactly \texttt{2026-08-22}: the fair \(n\)-category formulation is described as a proposed equivalence rather than a proved theorem in the cited report and current literature. This is supplied dated curation derived from the Oberwolfach material, not an independently verified literature assessment. Its status and cited claims require independent verification.
\end{itemize}
\end{document}
